\chapter{Stochastic Quantization And Singular SPDE}
\label{ch:constructive-stochastic-quantization-singular-spde}

\ConventionsEuclidean

This chapter introduces stochastic quantization as a construction and analysis
tool for Euclidean QFT.  The core object is a Markov process on fields whose
invariant distribution is the desired Euclidean measure or Schwinger
functional.  In dimensions where the equation is singular, renormalization is
part of the definition of the stochastic equation.

\section{Finite-Dimensional Langevin Identity}

Let \(S:\R^n\to\R\) be smooth and such that
\[
  Z=\int_{\R^n}\ee^{-S(x)}\dd x<\infty .
\]
The Langevin equation
\[
  \dd X_t=-\nabla S(X_t)\dd t+\sqrt2\,\dd B_t
\]
has generator
\[
  Lf=-\nabla S\cdot\nabla f+\Delta f .
\]
For smooth compactly supported \(f,g\),
\[
  \int (Lf)g\,\ee^{-S}\dd x
  =
  -\int \nabla f\cdot\nabla g\,\ee^{-S}\dd x .
\]
Therefore \(L\) is symmetric in \(L^2(\ee^{-S}\dd x)\), and the probability
measure \(Z^{-1}\ee^{-S}\dd x\) is invariant.  This identity is the finite
dimensional model for stochastic quantization.

\section{Regularized Field Equation}

Let \(C_\Lambda\) be a smooth ultraviolet cutoff covariance on a compact
Euclidean domain and let \(\xi\) be spacetime white noise in stochastic time
\(\tau\).  A regularized scalar \(\Phi^4\) Langevin equation has the form
\[
  \partial_\tau\phi_\Lambda
  =
  -C_\Lambda^{-1}\phi_\Lambda
  -\lambda_\Lambda \phi_\Lambda^3
  +c_\Lambda\phi_\Lambda
  +\xi_\Lambda .
\]
The noise covariance is chosen so that the linear part has invariant Gaussian
measure with covariance \(C_\Lambda\).  The drift is the \(L^2\)-gradient of
the regularized action
\[
  S_\Lambda(\phi)
  =
  \frac12(\phi,C_\Lambda^{-1}\phi)
  +
  \int\dd^dx\,
  \left[
    \frac{\lambda_\Lambda}{4}\phi^4
    -\frac{c_\Lambda}{2}\phi^2
  \right].
\]
The constants \(\lambda_\Lambda,c_\Lambda\) are coordinates of the
regularized theory.  They are tuned with \(\Lambda\) to produce a continuum
limit.

\section{Why Renormalized Products Enter}

For the Gaussian free field in dimension \(d\ge2\), \(\phi(x)\) is a
distribution rather than a pointwise random variable.  The square
\(\phi(x)^2\) is defined by a limiting procedure.  With mollifier
\(\rho_\epsilon\), set
\[
  \phi_\epsilon=\rho_\epsilon*\phi,\qquad
  :\!\phi_\epsilon^2\!:
  =
  \phi_\epsilon^2-\mathbb E[\phi_\epsilon^2].
\]
If the limit exists in a distribution space, it defines
\[
  :\!\phi^2\!: .
\]
Similarly,
\[
  :\!\phi_\epsilon^3\!:
  =
  \phi_\epsilon^3
  -3\mathbb E[\phi_\epsilon^2]\phi_\epsilon .
\]
Thus the renormalized stochastic equation for \(\Phi^4_3\) contains
counterterms because the nonlinear drift requires Wick powers, not pointwise
ordinary powers of the limiting field.

\section{Three-Dimensional Phi-Four Dynamic Equation}

On a three-dimensional torus, the renormalized dynamic equation is written
formally as
\[
  \partial_\tau\phi
  =
  \Delta\phi-m^2\phi-\lambda:\!\phi^3\!:+\xi .
\]
At cutoff \(\epsilon\), this is represented by
\[
  \partial_\tau\phi_\epsilon
  =
  \Delta\phi_\epsilon-m^2\phi_\epsilon
  -\lambda\phi_\epsilon^3
  +3\lambda C_\epsilon\phi_\epsilon
  +\xi_\epsilon
  +\text{finite coordinate terms},
\]
where
\[
  C_\epsilon=\mathbb E[\phi_\epsilon(x)^2]
\]
in the stationary Gaussian reference process.  The factor \(3\) is the number
of contractions in \(\phi^3\).  Additional finite coordinate terms depend on
the chosen renormalization scheme and on the desired mass coordinate.

\section{Invariant Measure And OS Data}

If the renormalized Markov process exists and is ergodic with invariant
measure \(\mu\), then Euclidean Schwinger functions are
\[
  S_n(f_1,\ldots,f_n)
  =
  \int \Phi(f_1)\cdots\Phi(f_n)\,\dd\mu(\Phi).
\]
To use these as QFT data one must verify:
\[
\begin{gathered}
  \text{Euclidean covariance},\qquad
  \text{reflection positivity},\\
  \text{regularity and growth bounds},\qquad
  \text{clustering or phase decomposition}.
\end{gathered}
\]
Stochastic construction alone gives a candidate invariant measure.  OS
reconstruction requires these additional properties.

\section{Relation To Constructive RG}

Stochastic quantization and constructive RG organize different maps.  The
stochastic equation gives a Markov semigroup in an auxiliary time
\(\tau\).  Constructive RG decomposes the covariance into scales and controls
effective interactions.  They meet when the stochastic dynamics is used to
construct the measure whose scale decomposition is then analyzed, or when RG
estimates prove convergence of the renormalized stochastic equation.

\begin{openproblem}[SPDE-to-Wightman construction]
\label{op:spde-to-wightman-construction}
Starting from a renormalized singular SPDE construction of an interacting
Euclidean field, prove the OS hypotheses, the linear growth condition needed
for reconstruction where applicable, the resulting Wightman distributions,
and the relation between stochastic counterterms and Wilsonian local
coordinates.
\end{openproblem}
